% SOMA — Paper 1: Selective Growth for Continual Learning
% Wakasa Labs · Nairobi, Kenya · March 2026
% Target: EMNLP 2026 or NeurIPS 2026

\documentclass[11pt,a4paper]{article}

% --- Packages ---
\usepackage[utf8]{inputenc}
\usepackage[T1]{fontenc}
\usepackage{amsmath,amssymb,amsfonts}
\usepackage{graphicx}
\usepackage{booktabs}
\usepackage{algorithm}
\usepackage{algorithmic}
\usepackage{hyperref}
\usepackage{xcolor}
\usepackage{natbib}

\title{%
    SOMA: Self-Organising Modular Architecture \\
    for Continual Learning via Selective Capacity Growth
}

\author{%
    Wakasa Labs \\
    Nairobi, Kenya \\
    \texttt{labs@wakasa.dev}
}

\date{March 2026}

\begin{document}

\maketitle

% =====================================================================
\begin{abstract}
SOMA (Self-Organising Modular Architecture) achieves Backward Transfer
BT $= -0.038$ on the Permuted MNIST continual learning benchmark using only
$K = 7$ LoRA adapters across 10 sequential tasks. This outperforms
Elastic Weight Consolidation (BT $\approx -0.09$), experience replay
(BT $\approx -0.07$), and unprotected sequential fine-tuning
(BT $\approx -0.18$). The core innovation is a three-signal necessity
detector that must unanimously agree before any new adapter capacity is
created, preventing the over-spawning failure mode that limits prior
modular approaches. An RL controller learns \emph{how} to grow ---
spawn, update, merge, or skip --- from a reward signal that penalises
forgetting twice as heavily as it rewards new-task accuracy. We validate
on both vision (Permuted MNIST, 10 tasks) and language (GSM8K, 3 subtasks
with Phi-3 Mini) benchmarks, demonstrating that structural prevention of
catastrophic forgetting is achievable without experience replay or
regularisation.
\end{abstract}

% =====================================================================
\section{Introduction}
\label{sec:intro}

% TODO: Write introduction
% Key points:
% - Catastrophic forgetting is the central problem of continual learning
% - Existing approaches: regularisation (EWC), replay, modular
% - Modular approaches spawn capacity but lack principled triggering
% - SOMA: necessity-gated growth with RL-guided action selection

% =====================================================================
\section{Related Work}
\label{sec:related}

% TODO: Review continual learning literature
% - EWC (Kirkpatrick et al., 2017)
% - GEM (Lopez-Paz & Ranzato, 2017)
% - KeepLoRA (2026), InfLoRA (CVPR 2024)
% - TreeLoRA (ICML 2025)
% - SplitLoRA

% =====================================================================
\section{Method}
\label{sec:method}

\subsection{Problem Definition}
% Definitions 2.1-2.4 from specification

\subsection{SOMA-NECESSITY: Three-Signal Capacity Trigger}
% N1, N2, N3 algorithms

\subsection{SOMA-GROW: RL-Guided Growth Controller}
% State, action, reward, REINFORCE

\subsection{SOMA-LEARN: Outer Training Loop}
% Cold start, main loop, metrics

% =====================================================================
\section{Experiments}
\label{sec:experiments}

\subsection{Experiment 1: Permuted MNIST}
% Table 1: SOMA vs baselines

\subsection{Experiment 2: GSM8K Sequential}
% Table 3: Language task results

% =====================================================================
\section{Ablation Study}
\label{sec:ablation}

% Table 2: Remove each Ni signal

% =====================================================================
\section{Conclusion}
\label{sec:conclusion}

% TODO: Summarise findings and future work (Papers 2, 3)

% =====================================================================
\bibliographystyle{plainnat}
\bibliography{references}

\end{document}
